% =====================================================================
\chapter{Inducción simple y fuerte; buen orden}
\label{tema:5}
% =====================================================================

Muchos enunciados del curso dependen de un número natural: la longitud de una secuencia, el número de capas de una red, el paso de un algoritmo iterativo. La inducción matemática permite demostrarlos para todos los valores a la vez, y es la técnica de demostración más usada en informática. Este tema presenta la inducción simple, la inducción fuerte y el principio del buen orden, y muestra que son tres formas de la misma idea. Se siguen el capítulo~10 de~\cite{hammack}, los capítulos~2 y~5 de~\cite{mcs} y las §6.1, 6.2 y 6.4 de~\cite{velleman}.

\section{El principio de inducción}

\begin{teorema}[Principio de inducción]\label{teo:induccion}
Sea $P(n)$ un predicado sobre $\N_0$. Si
\begin{enumerate}
  \item (\emph{caso base}) $P(0)$ es verdadera, y
  \item (\emph{paso de inducción}) para todo $n\in\N_0$, $P(n)\Rightarrow P(n+1)$,
\end{enumerate}
entonces $P(n)$ es verdadera para todo $n\in\N_0$.
\end{teorema}

La imagen habitual es la de una fila infinita de fichas de dominó: la primera cae, y cada una que cae derriba a la siguiente. En el paso de inducción, la suposición $P(n)$ se llama \emph{hipótesis de inducción}; no se está suponiendo lo que se quiere demostrar, sino probando un condicional para un $n$ arbitrario. Si el caso base es $P(n_0)$ con $n_0\in\Z$, la conclusión vale para todo $n\ge n_0$. El \cref{teo:induccion} se toma como axioma de los números naturales o se deduce del principio del buen orden (\cref{prop:bo-induccion}).

\addtocounter{definicion}{1}%

\addtocounter{definicion}{1}%

\addtocounter{definicion}{1}%

En las desigualdades, el paso de inducción suele necesitar una estimación intermedia, como aquí $2n+2\ge n+2$. Conviene escribir primero qué hay que obtener, $P(n+1)$, y buscar cómo llegar desde $P(n)$.

\addtocounter{definicion}{1}%

\section{Inducción fuerte}

\begin{teorema}[Inducción fuerte]\label{teo:induccion-fuerte}
Sea $P(n)$ un predicado sobre los enteros $n\ge n_0$. Si $P(n_0)$ es verdadera y, para todo $n\ge n_0$,
\[
\bigl(P(n_0)\wedge P(n_0+1)\wedge\dots\wedge P(n)\bigr)\Rightarrow P(n+1),
\]
entonces $P(n)$ es verdadera para todo $n\ge n_0$.
\end{teorema}

La inducción fuerte es útil cuando $P(n+1)$ depende no de $P(n)$, sino de casos anteriores más lejanos, como al descomponer un número en factores o en las recurrencias de varios términos.

\begin{teorema}[Existencia de la factorización en primos]\label{teo:factorizacion}
Todo entero $n\ge 2$ es primo o producto de primos.
\end{teorema}

\addtocounter{definicion}{1}%

\section{El principio del buen orden}

\begin{teorema}[Principio del buen orden]\label{teo:buen-orden}
Todo subconjunto no vacío de $\N_0$ tiene un elemento mínimo.
\end{teorema}

\begin{proposicion}\label{prop:bo-induccion}
El principio del buen orden implica el principio de inducción.
\end{proposicion}

La demostración anterior es el modelo de las \emph{pruebas por el menor contraejemplo}~\cite[§2.2]{mcs}, \cite[§10.3]{hammack}: para demostrar $\forall n\,P(n)$, se supone que el conjunto de contraejemplos no es vacío, se toma el menor y se construye un contraejemplo aún menor, o se llega a otra contradicción.

\begin{proposicion}[Existencia en la división euclídea]\label{prop:division-existencia}
Sean $a\in\Z$ y $b\in\N$. Existen $q,r\in\Z$ con $a=qb+r$ y $0\le r<b$.
\end{proposicion}

\section{Definiciones recursivas}

Los sumatorios, los productorios, las potencias y muchas sucesiones se definen \emph{recursivamente}: se da el primer valor y una regla para obtener cada valor a partir de los anteriores. Por ejemplo, $\sum_{k=1}^{0}a_k=0$ y $\sum_{k=1}^{n+1}a_k=\sum_{k=1}^{n}a_k+a_{n+1}$; la sucesión de Fibonacci se define por $F_1=F_2=1$ y $F_{n+2}=F_{n+1}+F_n$. La inducción es la forma natural de demostrar propiedades de objetos definidos así: el caso base comprueba la definición inicial y el paso usa la regla recursiva~\cite[§10.5]{hammack}.

\addtocounter{definicion}{1}%

\addtocounter{definicion}{1}%

